\exercisesection{3}

¶ 1. 设 A 是环，p 是 A 的素理想，B 是有限生成 A-代数，并设 \(r_{1} \subset \cdots \subset r_{m}\) 是 B 的位于 p 之上的非空递增理想序列。证明存在元素 \(a \in A - p\) 与 B 的元素的一个有限族 \((y_{k})_{1 \leqslant k \leqslant q}\)，具有下列性质：

\begin{enumerate}
\def\labelenumi{(\arabic{enumi})}
\item
  若记 \(\mathbf{B}' = \mathbf{A}[y_1, \ldots, y_q]\)，且 \(\mathbf{S}\) 是 \(a\) 的幂所成之集，则环 \(\mathbf{S}^{-1}\mathbf{B}\) 在 \(\mathbf{S}^{-1}\mathbf{B}'\) 上整。
\item
  对每个 \(i\)（\(1 \leqslant i \leqslant m\)），存在整数 \(h(i) \geqslant 0\)，使得 \(\mathfrak{r}_i\) 包含 \(y_1, \ldots, y_{h(i)}\)，且对每个多项式 \(\mathbf{P} \in \mathbf{A}[\mathbf{X}_{h(i)+1}, \ldots, \mathbf{X}_q]\)，关系 \(\mathrm{P}(y_{h(i)+1}, \ldots, y_q) \in \mathfrak{r}_i\) 蕴涵 \(\mathbf{P}\) 的所有系数都属于 \(\mathfrak{p}\)。
\end{enumerate}

此时理想 \(\mathfrak{r}_i\cap \mathbf{B}'\) 由 \(\mathfrak{p}\) 与诸 \(y_{k}\)（其中 \(k\leqslant h(i)\)）生成。（对 \(m\) 用归纳法，从而化归为 \(m = 1\) 的情形。然后设 \(\mathfrak{r} = \mathfrak{r}_1\alpha\)；对 A-代数 B 的生成元个数 \(n\) 用归纳法。设 \((x_{i})_{1\leqslant i\leqslant n}\) 是这样一个生成系，\(\mathfrak{s}\) 是 \(\mathrm{A}[\mathrm{X}_1,\dots ,\mathrm{X}_n]\) 中由满足 \(\mathrm{P}(x_1,\ldots ,x_n)\in \mathfrak{r}\) 的那些 P 组成的理想；则 \(\mathfrak{s}\supset \mathfrak{pA}[\mathrm{X}_1,\ldots ,\mathrm{X}_n]\)；按 \(\mathfrak{s}\) 是否等于 \(\mathfrak{pA}[\mathrm{X}_1,\ldots ,\mathrm{X}_n]\) 区分两种情形。在第二种情形，存在多项式 \(\mathrm{Q}\in \mathrm{A}[\mathrm{X}_1,\ldots ,\mathrm{X}_n]\)，它的所有系数都不在 \(\mathfrak{p}\) 中，且它属于 \(\mathfrak{s}\)。证明可以确定整数 \(m_i\)（\(2\leqslant i\leqslant n\)），使得当记 \(x_1' = \mathrm{Q}(x_1,\ldots ,x_n),x_i' = x_i - x_1^{m_i}\)（\(2\leqslant i\leqslant n\)）时，存在元素 \(b\in \mathrm{A} - \mathfrak{p}\)，使 \(bx_{1}\) 在 \(\mathrm{C} = \mathrm{A}[x_2',\ldots ,x_n']\) 上整；把归纳假设应用于 A-代数

\[
\mathrm{B} _ {1} = \mathrm{A} [ x _ {2} ^ {\prime}, \dots , x _ {n} ^ {\prime} ]
\]

以及理想 \(\mathfrak{r} \cap \mathbf{B}_1\)。）

\begin{enumerate}
\def\labelenumi{\arabic{enumi}.}
\setcounter{enumi}{1}
\item
  设 K 是域，B 是两个未定元的多项式环 K{[}X, Y{]} 对 K{[}X, Y{]} 的由 \(X^{2}\) 与 XY 生成的理想 a 的商环；用 x、y 表示元素 X、Y 在 B 中的像；设 A 是 B 的子环 K{[}x{]}。证明元素 x 与 \(y^{n}\)（\(n \geqslant 0\)）构成 B 在 K 上的一组基，且 B 不在 A 上整。另一方面，B 中不存在在 A 上代数自由的元素。由此给出当 A 不是整环时第 1 节定理 1 的推论 1 的一个反例。
\item
  设 K 是无限域，L 是 K 的扩域，\((x_{i})_{1 \leq i \leq n}\) 是 L 中元素的有限族，d 是 \(\mathrm{K}(x_{1}, \ldots, x_{n})\) 在 K 上的超越次数。证明存在 L 的 d 个元素 \(z_{j} = \sum_{i=1}^{n} a_{ji} x_{i}\)（其中 \(a_{ji} \in K\)），使得 \(\mathrm{K}[x_{1}, \ldots, x_{n}]\) 在 \(\mathrm{K}[z_{1}, \ldots, z_{d}]\) 上整。若进一步设 \(\mathrm{K}(x_{1}, \ldots, x_{n})\) 是 K 的可分扩张，证明诸 \(a_{ji}\) 可如此选取，使诸 \(z_{j}\) 构成 \(\mathrm{K}(x_{1}, \ldots, x_{n})\) 在 K 上的一个分离超越基（当 n \textgreater{} d 时对 n 用归纳法）。
\item
  设 B 是整环，A 是 B 的子环。设存在 B 的 \(n \geqslant 1\) 个元素 \(t_{1}, \ldots, t_{n}\)，它们在 A 上代数无关，且 B 在 \(A[t_{1}, \ldots, t_{n}]\) 上整。设 n 是 B 的一个极大理想，\(p = n \cap A\)。则存在 B 的素理想 q，它含于 n、异于 n 且包含 p。（先通过考虑 A 的整闭包 \(A'\)（含于 B 的分式域中）与环 \(B' = B[A']\)，化归为 A 整闭的情形；在 A 整闭的情形，利用 §2 第 4 节定理 3。）
\end{enumerate}

¶ 5. 设 A 是整环，B ⊃ A 是 A 的分式域 K 的子环；设 A 在 B 中整闭，且 B 包含元素 t，使 B 是有限生成 A{[}t{]}-模。

\begin{enumerate}
\def\labelenumi{(\alph{enumi})}
\tightlist
\item
  设 F 是一个 \(\neq 0\) 的多项式，属于 A{[}X{]}，使 \(F(t)\) 属于导子 \(\mathfrak{f}\)（B 相对于 A{[}t{]} 的）；证明：若 \(c\) 是 F 的首项系数，则 \(c\) 属于 \(\mathfrak{f}\) 的根。（通过考虑环 A{[} \(c^{-1}\) {]} 与 B{[} \(c^{-1}\) {]}，化归为证明：当 \(c = 1\) 时必有 B = A{[}t{]}。另一方面，可设 F(t) \(\neq 0\)，否则 \(t \in \mathbf{A}\)。对每个 \(x \in \mathbf{B}\)，由假设
\end{enumerate}

\[
x \mathrm{F} (t) = \mathrm{F} (t) \mathrm{Q} (t) + \mathrm{R} (t)
\]

其中 Q 与 R 是 A{[}X{]} 中的多项式，且 \(\deg(\mathrm{R}) < \deg(\mathrm{F})\)；只需考虑 \(\mathrm{R}(t) \neq 0\) 的情形；若 \(y = \mathrm{R}(t)(\mathrm{F}(t))^{-1} = x - \mathrm{Q}(t)\)，证明 t 在 \(A[y^{-1}]\) 上整，并由此推出 y 在 A 上整。）

\begin{enumerate}
\def\labelenumi{(\alph{enumi})}
\setcounter{enumi}{1}
\item
  再设 A 是 Noether 的；设 q 是属于 \(\operatorname{Ass}(B/\mathfrak{f})\) 的 B 的素理想。证明：若 \(\omega\) 是从 B 到 \(B/q\) 的分式域的典范同态，则 \(\omega(t)\) 在 \(\mathbf{A}/(\mathfrak{q} \cap \mathbf{A})\) 的分式域上是超越的。（若 \(p = q \cap A[t]\)，则存在元素 \(y \in B\)，使 q 是由那些 \(z \in B\)（满足 \(zy \in \mathfrak{f}\)）组成之集；证明 p 是 \(B' = A[t] + By\) 相对于 A{[}t{]} 的导子。然后利用 (a) 以反证法论证。）
\item
  设 A 是 Noether 的。设 n 是 B 的一个极大理想，\(p = A \cap n\)。证明：若 \(A_{p} \neq B_{n}\)，则存在 B 的素理想，它含于 n、异于 n 且包含 p。（设 \(\mathfrak{f}\) 是 B 相对于 A{[}t{]} 的导子。先设 \(n \supset \mathfrak{f}\)；在 B 的含于 n 且包含 \(\mathfrak{f}\) 的素理想所成之集中考虑一个极小元 q；考虑 A 在 B/q 中的典范像，并利用 (b) 与习题 4。若 \(n \not\supset \mathfrak{f}\)，则局部环 B\_n 等于局部环 A{[}t{]}\_\{n\_1\}，其中 n\_1 = n ∩ A{[}t{]}（§ 1 习题 16）。若 t ∈ A\_p，证明 A\_p = B\_n。反之，若 A\_p ≠ B\_n，则由 § 2 习题 14 (c) 推出 pA\_p{[}t{]} 是 A\_p{[}t{]} 的一个非极大素理想，并断定 pB\_n 是 B\_n 的异于 nB\_n 的素理想。）
\end{enumerate}

\begin{enumerate}
\def\labelenumi{\arabic{enumi}.}
\setcounter{enumi}{5}
\tightlist
\item
  若整环 A 是 Noether 的，且对 A 的分式域的扩域中元素的每个有限序列 \((t_{i})_{1\leq i\leq n}\)，\(A[t_{1},\ldots,t_{n}]\) 的整闭包都是有限生成 \(A[t_{1},\ldots,t_{n}]\) -模，则称 A 为整 Noether 整环。域上的每个有限生成整代数都是整 Noether 整环（第 2 节定理 2）。
\end{enumerate}

\begin{enumerate}
\def\labelenumi{(\alph{enumi})}
\item
  整 Noether 整环上的每个有限生成整代数都是整 Noether 整环。不约化为 0 的整 Noether 整环的每个分式环都是整 Noether 整环。
\item
  设 A 是整 Noether 整环，\((t_{i})_{1\leq i\leq n}\) 是 A 的分式域 K 的扩域中元素的一个有限序列，B 是 \(\mathbf{K}(t_{1},\ldots,t_{n})\) 的子环，包含 A 且在 A 上整。则 B 是有限生成 A-模。（利用《代数》第五章 § 5 第 7 节，注意 B 的分式域是 K 的有限代数扩张。）
\end{enumerate}

\begin{enumerate}
\def\labelenumi{\arabic{enumi}.}
\setcounter{enumi}{6}
\item
  设 A 是整 Noether 整环（习题 6），\((t_i)_{1 \leq i \leq n}\) 是 A 的分式域的扩域中元素的一个有限序列，B 是环 \(A[t_1, \ldots, t_n]\)，\(n\) 是 B 的一个极大理想，\(p = A \cap n\)，\(A'\) 是 A 在 B 中的整闭包，\(p' = n \cap A'\)。那么，若局部环 \(A_p'\) 与 \(B_n\) 不同，则存在 B 的一个素理想，含于 \(n\)、异于 \(n\) 且包含 \(p\)（“Zariski 主定理”）。（利用习题 6 (b)，化归为 \(A' = A\) 的情形，再化归为 A 是以 \(p\) 为极大理想的局部环的情形。证明此时 \(B/n\) 是 \(A/p\) 的有限代数扩张（参见第 4 节定理 3）；由此推出，对每个满足 \(A \subset C \subset B\) 的环 C，\(n \cap C\) 都是 C 的极大理想。然后对 \(n\) 用归纳法，并用 \(B_t\) 表示 \(A[t_1, \ldots, t_i]\) 在 B 中的整闭包；按 \(t_{i+1}\) 在 \(B_t\) 的分式域上是超越的还是代数的区分两种情形；在前一情形用习题 4，在后一情形用习题 5 (c)。）
\item
  设 A 是环。为使 A 是 Jacobson 环，必须且只需：对每个极大理想 \(m'\)（多项式环 A{[}X{]} 的），\(m' \cap A\) 都是 A 的极大理想。（注意：若整环 B 有 Jacobson 根 \(R \neq 0\)，且 \(b \in R\)，\(b \neq 0\)，则 B 与 B{[}X{]} 的包含 \(1 + bX\) 的一个极大理想之交不能包含 b。）
\item
  给出一个 Jacobson 整环 A 与包含 A 的 Jacobson 环 B 的例子，使得存在 B 的一个极大理想，它与 A 之交不是极大理想。
\end{enumerate}

¶ 10. 设 k 是域，L 是无限集，A 是多项式环 \(k[X_{\lambda}]_{\lambda \in L}\)。若 \(k_{0}\) 是 k 中所含的素域，用 b 表示 k 在 \(k_{0}\) 上的一个超越基的基数，并记 \(c = \text{Card}(L)\)。

\begin{enumerate}
\def\labelenumi{(\alph{enumi})}
\item
  对 A 的每个理想 a，证明存在 a 的一个生成系 S，使 \(\operatorname{Card}(S) \leqslant c\)。（对 L 的有限子集 J，考虑 a 与子环 \(k[X_{\lambda}]_{\lambda \in J}\) 的交。）
\item
  设 \(c < b\)。证明：若 \(m\) 是 \(A\) 的一个极大理想，则 \(A / m\) 是 \(k\) 的代数扩张。（设 \(K\) 是 \(k\) 的在 \(k_0\) 上由 \(m\) 的一个基数 \(\leqslant c\) 的生成系中诸多项式的系数生成的子域；设 \(m_0 = m \cap K[X_\lambda]_{\lambda \in L}\)，于是 \(m = m_0 \otimes_k k\)，且
\end{enumerate}

\[
\mathrm{A} / \mathfrak {m} = (\mathrm{K} [ \mathrm{X} _ {\lambda} ] _ {\lambda \in \mathrm{L}} / \mathfrak {m} _ {0}) \otimes_ {\mathrm{K}} k.
\]

最后注意，存在从 \(K[X_{\lambda}]_{\lambda \in L}/m_{0}\) 的分式域到 k 的一个代数闭包上的 K-同构。）由此推出，此时 A 是 Jacobson 环（利用习题 8）。

\begin{enumerate}
\def\labelenumi{(\alph{enumi})}
\setcounter{enumi}{2}
\tightlist
\item
  设 \(b < c\)；则 \(\operatorname{Card}(k) \leqslant c\)，从而存在双射 \(\lambda \mapsto \xi_{\lambda}\)，把一个子集 \(L'\)（\(L\) 的）映到 \(k\) 上，且 \(L'\) 在 \(L\) 中的余集非空。设 \(\lambda_0 \in L - L'\)；证明存在一个 \(k\) -同态 \(u\)，它把 \(A\) 映到有理函数域 \(k(Y)\) 的一个子域上，使得 \(u(X_{\lambda_0}) = Y\)，
\end{enumerate}

\[
u (\mathbf {X} _ {\lambda}) = 1 / (\mathbf {Y} - \xi_ {\lambda})
\]

（\(\lambda \in \mathbf{L}'\)），且 \(u(\mathbf{X}_{\lambda}) = 0\)（对 \(\lambda \in \mathbf{L} - (\mathbf{L}' \cup \{\lambda_0\})\)）；若 \(\mathfrak{m}\) 是 \(u\) 的核，则 \(\mathbf{A} / \mathfrak{m}\) 因此不是 \(k\) 的代数扩张。

\begin{enumerate}
\def\labelenumi{(\alph{enumi})}
\setcounter{enumi}{3}
\tightlist
\item
  在与 (c) 相同的假设下，证明 A 不是 Jacobson 环。（注意存在不是 Jacobson 环的 k-代数 B（例如局部环），使 \(\operatorname{Card}(B) = b\)。）
\end{enumerate}


